\documentclass[11pt]{article}
% === core_packages.tex ===
\usepackage{amsmath, amssymb, amsthm, amsfonts}
\usepackage{geometry}
\usepackage{hyperref}
\usepackage{graphicx}
\usepackage{physics}
\usepackage{booktabs}
\usepackage{listings}
\usepackage{bookmark}
\usepackage{amscd}
\usepackage{tikz-cd}
\usepackage{mathtools}
\usepackage{xspace}
\usepackage[backend=biber, style=numeric]{biblatex}
\usepackage{tcolorbox}
\tcbuselibrary{theorems}
\usepackage{parskip}
\usepackage{tikz}
\usepackage{etoolbox}
\usetikzlibrary{arrows.meta, shapes, positioning, calc, cd, shapes.geometric, backgrounds}
\usepackage[en-US]{datetime2}
\usepackage{array}


% --- Math & Ontology Macros ---
\newcommand{\entropy}{S}
\newcommand{\opponent}{opponent}
\newcommand{\aside}[1]{[\emph{\opponent}: ``#1'']}
\newcommand{\paper}[1]{[\emph{#1}: ``#1'']}
\newcommand{\iame}{Real\textsuperscript{e}}
\newcommand{\iameverything}{Real\textsuperscript{everything}}
\newcommand{\fabric}{voxel}
\newcommand{\Fabric}{Voxel}
\newcommand{\pdflink}[1]{}



% --- Global Layout Defaults ---
\geometry{letterpaper, margin=1in}
\hypersetup{
    colorlinks=true,
    linkcolor=black,
    citecolor=black,
    filecolor=magenta,
    urlcolor=blue
}
\hfuzz=5pt \vfuzz=5pt \hbadness=10000 \vbadness=10000

\author{\iame}


\usepackage{tocloft}

% --- Article Specific Theorems ---
\newtheorem{lemma}{Lemma}
\theoremstyle{definition}
\newtheorem{axiom}{Axiom}
\newtheorem{theorem}{Theorem}
\newtheorem{corollary}{Corollary}
\newtheorem{definition}{Definition}
\newtheorem{proposition}{Proposition}
\newtheorem{conjecture}{Conjecture}


\makeatletter
\let\remark\@undefined
\let\endremark\@undefined
\makeatother
\theoremstyle{remark}
\newtheorem*{remark}{Remark}

\newtcbtheorem[number within=section]{principle}{Principle}%
{colback=blue!5,colframe=blue!50!black,fonttitle=\bfseries}{principle}

% --- Article TOC Layout ---
\setlength{\cftbeforesecskip}{1.0em}
\renewcommand{\cftsecleader}{\cftdotfill{\cftdotsep}}
\renewcommand{\cftsubsecleader}{\cftdotfill{\cftdotsep}}


\addbibresource{../references.bib}

\title{Emergence of Quantum Mechanics: Fermions, Bosons and Pauli Exclusion}
\author{Juha Meskanen}
\date{June 2026}

\begin{document}

\maketitle

\begin{abstract}
  Papers~V \cite{meskanen2026v} established that Quantum Mechanical wavefunction
  is the compression codec.
\end{abstract}

\tableofcontents
\newpage

% ============================================================
\section{Introduction}
\label{sec:intro}

Bosons don't exist as fundamental objects. They are explanatory fictions — the story we tell
about why fermions (pixels) move the way they do, when we don't know
they're being sampled from a wavefunction.

The actual causal chain is:

\[
\text{wavefunction (codec)} \xrightarrow{\text{Born rule sampling}} \text{fermion configurations (pixel frames)}
\]

But an observer who only sees the pixel outputs, not the codec, invents a story:

\[
\text{fermion} \xrightarrow{\text{"virtual boson exchange"}} \text{fermion}
\]

The boson is the compression residual made into a noun.
It's not wrong — it's a valid effective description — but it's not fundamental.

The analogy is exact: In MPEG, macroblocks don't "exchange" anything.
The DCT coefficients just are what they are.
But if one only watched the decoded pixels and tried to explain their correlations
without knowing about DCT, one would invent something that looks exactly like boson exchange.



% ============================================================
\section{Setup}
\label{sec:setup}

We take a sequence of pixel frames. We compress them with $\psi$.
We hypothetize that the off-diagonal residuals — the "bosons" — reproduce the correct correlation structure
between fermion positions, without any interaction term being put in by hand.


\section{2x2 Minimal Video}

TODO

The minimum complexity states ($C_s = 0$) — uniform frames, all pixels same —
already carry maximum boson signal $\sqrt{3}/2$. And crucially, videos $5 \rightarrow 10$ which
are checkerboard patterns flipping to their inverse, also hit $\sqrt{3}/2$ at $C_s = 3$.

The vacuum configuration (uniform, zero information content) and the maximally
anti-correlated configuration (checkerboard flip) produce identical boson signal.

That's not trivial — it suggests the boson signal is tracking something topological about
the pixel transition, not just the complexity.


\section{Experiment 2:}

Instead of just measuring the norm of the off-diagonal, we look at the structure of the
off-diagonal matrix — specifically:
\begin{itemize}
\item Is it symmetric or antisymmetric? (spin-0 vs spin-1 candidate)
\item Does it propagate? (is there a spatial pattern that shifts between frames?)
\item Is it quantised? (does the signal come in discrete units?)
\end{itemize}



\subsection*{The key result — video (1,2):}

A single fermion hops from site 0 to site 1. The boson matrix is purely antisymmetric — not approximately,
exactly. Its eigenvalues are $\pm 0.5$, symmetric about zero.
And the forward hop $(1,2)$ and reverse hop $(2,1)$ produce identical boson structure
with opposite momentum $k = \pm 0.157$.

That is exactly what one would expect from a particle/antiparticle pair.

\textbf{The block hop (3,12):}

Two fermions move together $[1100] \rightarrow [0011]$. The boson has momentum $k = 2\pi/10$ —
a longer wavelength, lower energy mediator carrying two fermions.
Different particle, different $k$.

What this seems to demonstate: the boson isn't just a correlation residual in the abstract. It has:
\begin{itemize}
\item Definite antisymmetry — fermionic exchange signature built in
\item Definite momentum — propagation direction from the hop direction  
\item Particle/antiparticle pairing — forward and reverse hops are related by $k \rightarrow -k$
\end{itemize}

All from the codec structure alone. No interaction term put in by hand.

The claim: the compression residual of a single fermion hop is a purely antisymmetric,
propagating, quantised object with opposite momentum for the reverse process.
That looks like a photon.

\section{Experiment 3: More Complex Videos}

We scale the video up — longer chains, more frames — to see if the momentum
quantisation follows $k = 2\pi n / L$?

\subsection*{The key result — analytically exact:}

For any single fermion hop $0 \rightarrow d$, regardless of $L$ or $d$:

\[
\psi = \frac{1}{\sqrt{2}}(e_0 + i \cdot e_d)
\]

The boson matrix has exactly two non-zero entries:

\[
B_{0,d} = +\frac{i}{2}, \quad B_{d,0} = -\frac{i}{2}
\]

Always purely antisymmetric. Always eigenvalues $\pm\frac{1}{2}$. Always.

And the phase of $B_{0,d}$ is always $-\pi/2$, so:

\[
k \cdot d = -\frac{\pi}{2} \quad \Rightarrow \quad k = -\frac{\pi}{2d}
\]

This isn't a momentum in the plane-wave sense — it's a phase relationship between the two fermion sites.
The boson encodes exactly one number: the $\pi/2$ phase shift between where the fermion was and where it went.

Physically the boson is a quarter-turn in phase space. That's a propagator. That's exactly what a photon does in QED.

The real structure is $\pi/2$ invariant.

The compression residual of a single fermion hop is an antisymmetric, purely imaginary,
$\pm\frac{1}{2}$-eigenvalued object carrying a universal $\pi/2$ phase — independent
of chain length and hop distance.


\section{Experiment 4}

We check whether two-fermion hops break this universality, and whether that breaking has physical meaning.

\subsection*{Result 1: Universality of Phase}

The $\pi/2$ phase is universal for single hops — and partially survives in two-fermion hops.

Every two-fermion boson matrix contains entries with phase $-\pi/2$.
But it also acquires new phases: $0$ and $+\pi/2$.
The boson matrix for two fermions is a superposition of single-hop bosons — one per fermion pair.
The $-\pi/2$ entries are the individual hops. The $0$ and $+\pi/2$ entries are the inter-fermion correlations
— the part that has no single-particle analogue.

\subsection*{Result 2: Path Independence}

Crossing is indistinguishable from non-crossing.
Fermions passing through each other produce the identical boson matrix as
fermions going the same direction.
The codec doesn't care about trajectory — only about initial and final pixel configurations.
This is path independence falling out for free.


\subsection*{Result 3: Pauli Exclusion}

The Pauli test is the most striking result:

\begin{itemize}

\item Two fermions at the same site: $\|\psi\| = \sqrt{2}$, boson matrix is exactly zero.
  No compression residual. No boson.
  The wavefunction is real-valued and uniform — the codec has nothing to encode.
  Double occupancy is not forbidden by a rule — it produces no boson, meaning it's
  invisible to the codec. It simply doesn't compress.

\item Two fermions hopping to the same site: boson matrix is non-zero but the phases
  break the $\pi/2$ universality.
  The codec is struggling — it's trying to describe something the pixel model can't cleanly represent.

\end{itemize}

Single-fermion hops $\rightarrow$ universal $-\pi/2$ boson (the photon).

Double occupancy $\rightarrow$ zero boson (Pauli exclusion = zero compressibility).

Multi-fermion hops $\rightarrow$ superposition of single bosons + interaction residuals.

Pauli exclusion isn't a postulate — it's the statement that double occupancy has zero compression residual,
making it informationally invisible.




% ============================================================
\section{Discussion}
\label{sec:discussion}


% ============================================================
\section{Future Work}
\label{sec:future}



% ============================================================
\section*{Supplementary Material}
\label{sec:supplementary}



\printbibliography

\end{document}
